\section{Direct quadratic batching and a nine-variable equation}
\label{sec:nine}

\subsection{The exact positive quadratic source}

We start with the literal 31-variable source of package~30, retained
by the later report packages~\cite{repo}. Its strictly positive unknowns,
in fixed order, are
\[
\begin{gathered}
 J,F,\alpha,z,f,h,i,o,j,s,w,\tau,\eta,\zeta,y,\gamma,\\
 N,E,P,c,a,X,D,C_2,L,F_2,V,g,I,B_a,\Delta.
\end{gathered}
\]
Only in this subsection, put $q=\beta J+1$, $k=\eta+\zeta$,
$H=4a+3$, and $C_w=q-F-\alpha$. The fixed ports are
$(\beta,K_0,M_C,M_F)$ from Section~\ref{sec:complemented-compiler}.
The 22 residuals, each required to vanish, are
\begin{alignat}{2}
 &L-q^2+1, &&N-q(L+1),\nonumber\\
 &X-wq, &&a-E-sN,\nonumber\\
 &E-X(a-E), &&c-k(a-E)-\eta,\nonumber\\
 &P-E(c-\eta)-k, &&\tau^2-P(P-k)-1,\nonumber\\
 &D-X-ac-\gamma H, &&\Delta-a^2-H,\nonumber\\
 &C_2-c^2, &&D^2-\Delta C_2-1,\nonumber\\
 &k-jc+V-1-hE, &&\label{eq:parent22}
\end{alignat}
\vspace{-6mm}
\begin{align*}
 &jc-V-(q-F)(N-q)+C_wL-(M_C+qM_F)J,\\
 &(K_0+w)C_w+q-F-z(q-1)-1,\\
 &F_2-f^2,\qquad V-of+c,\qquad g-y^2+V^2,\\
 &I-iC_2,\qquad B_a-I^2,\qquad B_a-\Delta(F_2-1),\\
 &y^2-B_ag-1.
\end{align*}
Let $r_0,\ldots,r_{21}\in\Z[n_1,\ldots,n_{31}]$ be these residuals
after replacing every positive source coordinate by $n_i+1$.
The delivered literal parent file \texttt{nine/parent31.json} has SHA-256
\begin{center}\footnotesize\ttfamily
1c336ba0b3c9827ffddcf7ce895cad9dca2907f907cde7484c7f387bb2cb1287.
\end{center}
Its historical binding metadata is preserved as provenance; the active
binding of this report is the new complete sibling compiler tuple.

For clarity, the elimination producing the 31-variable source is exact
on the positive domain. Its former register was
$G_p=(q-1)(q-F)+\alpha=q(q-F)-C_w$.
On $N=q(L+1)$ one has
\[
 LG_p=(q-F)(N-q)-C_wL.
\]
The transport row gives
$(K_0+w)C_w+(q-F)=1+z(q-1)\ge q$.
If $C_w\le0$, its left side is at most $q-F<q$, a contradiction.
Thus $C_w>0$, $q-F>0$, and the reconstructed $G_p$ is positive.
The unchanged normalized compiler theorem therefore gives
\begin{equation}
 \exists\mathbf n\in\N^{31}\ \bigwedge_{j=0}^{21}r_j(\mathbf n)=0
 \quad\longleftrightarrow\quad
 \text{K215 halts on blank input}.
 \label{eq:parent-calibration}
\end{equation}
The fifth residual at $\mathbf n=0$ is $1$, so a simultaneous zero
is never the all-zero vector. We use this small fact in completeness.

\subsection{A finite Sidon certificate replaces exponential positions}

Use the following increasing list of 32 positions:
\begin{equation}
\begin{split}
 (0,a_1,\ldots,a_{31})={}&(0,10,35,55,62,105,144,145,159,168,232,245,\\
 &306,323,325,353,371,374,382,494,499,525,566,578,\\
 &610,647,681,685,741,777,783,799).
\end{split}
\label{eq:weights}
\end{equation}
All 528 unordered pair sums with repetition are distinct. The file
\texttt{weights.json} and the independent verifier give a complete
finite certificate. No unverified extremality claim for this set is
needed. Put
\[
 M=799,\qquad d=4M+1=3197,\qquad \mathcal H=22d=70334.
\]
Write
\[
 r_j(\mathbf n)=\sum_{|\alpha|\le2}p_{j,\alpha}\mathbf n^\alpha,
 \qquad
 d_{j,\alpha}=(2-|\alpha|)!\prod_{i=1}^{31}\alpha_i!\,p_{j,\alpha}.
\]
Define the following fixed polynomials and fixed target positions:
\begin{align}
 E_j(T)&=\sum_{|\alpha|\le2}d_{j,\alpha}
                    T^{2M-\sum_i a_i\alpha_i},\nonumber\\
 E_*(T)&=\sum_{j=0}^{21}T^{jd}E_j(T),&
 Z_*(T)&=\sum_{j=0}^{21}T^{jd+2M}.
 \label{eq:kernel}
\end{align}
There are 136 nonzero weighted coefficient records and only 16
distinct nonliteral coefficient recipes. They are all explicit
polynomials in the four fixed native ports.

\begin{lemma}[Simultaneous quadratic extraction]
\label{lem:quadratic-extraction}
For $c(T)=1+\sum_{i=1}^{31}n_iT^{a_i}$,
\[
 [T^{jd+2M}]\,c(T)^2E_*(T)=2r_j(\mathbf n)
 \qquad(0\le j<22).
\]
\end{lemma}
\begin{proof}
Pad a monomial of total degree at most two with zero positions to
obtain a pair from the set in \eqref{eq:weights}. The pair-sum
certificate makes this representation unique. The coefficient of its
monomial in $c(T)^2$ is
$2!/((2-|\alpha|)!\prod_i\alpha_i!)$; this cancels the factorial
weight in $d_{j,\alpha}$, leaving $2p_{j,\alpha}\mathbf n^\alpha$.
Finally $c(T)^2E_j(T)$ has support in $[0,4M]$. Translating it by
$jd$ puts it in $[jd,jd+4M]$; these intervals are disjoint.
Thus no other row contributes to a target.
\end{proof}

The important distinction is between a polynomial's degree and the
choice of a positional code for its monomials. A quadratic system
needs unique pair sums, not separate base-five positions for a quartic
in 31 coordinates. Moreover the 22 target tests can share the same
unknown digit code; one does not need an extra unknown for each row.

\subsection{One explicit size constant}

Expand each weighted coefficient into monomials in the four positive
native ports. Let $C_*$ be one plus the sum of the absolute values of
all scalar coefficients times their corresponding port monomials,
over all 136 records. Thus every $|d_{j,\alpha}|<C_*$.
This is a bounded, finite integer expression, independent of the
existence of any witness. Set
\begin{equation}
 \kappa=2^{\bitlength(2048C_*)}>2048C_*.
 \label{eq:kappa}
\end{equation}
For the delivered tuple,
\[
 \bitlength(C_*)=1\,359\,514,\qquad
 \boxed{\kappa=2^{1\,359\,525}.}
\]
The constant's unsigned little-endian binary magnitude has 169,941
bytes; its compressed payload and all recipe hashes are included.
This replaces any unspecified coefficient-height parameter. The
derived fixed coefficients, including negative ones, remain signed
integer leaves of the final circuit.

\subsection{The new source code and its two-way interpretation}

For the moment let $b\ge2$ and $g>0$ be integers. Set
\begin{align}
 B&=\kappa b^2,&Q_0&=B^{M+1},&
 \ell&=\sum_{i=1}^{31}B^{a_i},\nonumber\\
 \lambda&=\sum_{t=0}^{\mathcal H-1}B^t,&
 S_0&=2(1+g)^2E_*(B)+B\lambda,\nonumber\\
 S&=g+Q_0S_0,&
 T&=Q_0-1-(b-1)\ell+Q_0(B-2)Z_*(B),\nonumber\\
 N&=B^{\mathcal H+M+2}=B^{71135},&
 R&=S(N^2-N)+(T+1)(N^2-1),\nonumber\\
 G_{\rm cap}&=32bB^M.&&
 \label{eq:nine-source}
\end{align}
Everything in this display is an integer polynomial in $b,g$ with
fixed integer coefficients. In particular the finite sum defining
$\lambda$ introduces no division in the polynomial.

\begin{lemma}[Unconditional nonvanishing and conditional size]
\label{lem:source-size}
For every $b\ge2,g>0$, the integer $R$ is nonzero. If additionally
$g<4bB^M$, then $0\le S,T<N$ and $R\ge N^2-1\ge N\ge\max\{b,8\}$.
\end{lemma}
\begin{proof}
Since $Q_0\mid N$, reduction of \eqref{eq:nine-source} modulo $Q_0$
gives
\[
 R\equiv (b-1)\ell\pmod{Q_0}.
\]
The positions in $\ell$ are distinct and at most $M$, while $b<B$.
Consequently $0<(b-1)\ell<Q_0$, proving $R\ne0$ without a sign
assumption on $E_*(B)$ or $R$.

Now suppose $g<4bB^M$. We have $g<Q_0$, because $B>4b$.
The highest exponent in $E_*$ is at most $\mathcal H-2M-1$.
Using $B\ge2$ and the coefficient bound,
\[
 |E_*(B)|<2C_*B^{\mathcal H-2M-1},\qquad
 1+g<5bB^M.
\]
Thus
\[
 |2(1+g)^2E_*(B)|
 <100C_*b^2B^{\mathcal H-1}<B^{\mathcal H}.
\]
Since $B^{\mathcal H}\le B\lambda<2B^{\mathcal H}$, it follows
that $0<S_0<3B^{\mathcal H}$. The low mask
$Q_0-1-(b-1)\ell$ is between $0$ and $Q_0-1$: its base-$B$ digits
are $B-b$ at the chosen positions and $B-1$ elsewhere. Also
$(B-2)Z_*(B)<B^{\mathcal H}$. The extra factor of $B$ in $N$
therefore gives $0\le S,T<N$. Finally $T+1\ge1$ and $S\ge0$ imply
$R\ge N^2-1$; the remaining bounds are immediate.
\end{proof}

Write $\tau_2(u,v)=0$ for addition without a binary carry, equivalently
disjoint sets of occupied binary positions for nonnegative $u,v$.

\begin{lemma}[All residuals in one carry test]
\label{lem:carry-decoder}
Suppose $b$ is a power of two, $b\ge2$, and $0<g<4bB^M$. Then
\[
 \tau_2(S,T)=0
 \quad\longleftrightarrow\quad
 \begin{cases}
 g=\sum_i n_iB^{a_i},\quad 0\le n_i<b,\\
 r_0(\mathbf n)=\cdots=r_{21}(\mathbf n)=0.
 \end{cases}
\]
Conversely, every simultaneous natural zero can be encoded in this
way for arbitrarily large powers of two $b$, with $0<g<bB^M$.
\end{lemma}
\begin{proof}
Now $B$ and $Q_0$ are powers of two. The strict block bounds in
Lemma~\ref{lem:source-size} split the carry condition into
\[
 \tau_2\bigl(g,Q_0-1-(b-1)\ell\bigr)=0,
 \qquad
 \tau_2\bigl(S_0,(B-2)Z_*(B)\bigr)=0.
\]
The first mask has all bits set at unchosen digit positions. At a
chosen position its digit $B-b$ has precisely the bits above the
lowest $\log_2b$ bits set. Hence it gives exactly the stated code
and digit bounds.

For decoded digits let $h_t$ be the coefficient of $T^t$ in
$c(T)^2E_*(T)$. Its degree is at most $\mathcal H-1$, and
\[
 |h_t|\le C_*\left(1+\sum_i n_i\right)^2
 <1024C_*b^2<B/2.
\]
Therefore the integer
\[
 S_0/2=(1+g)^2E_*(B)+(B/2)\lambda
\]
has ordinary base-$B$ digits $B/2+h_t$, all strictly between $0$
and $B$. Both summands in the second carry test are even. Divide
them by two; the mask becomes $(B/2-1)Z_*(B)$.
At each chosen target digit, $B/2-1$ occupies all the lower binary
positions. A digit strictly between $0$ and $B$ is disjoint from
this mask exactly when it equals $B/2$. By
Lemma~\ref{lem:quadratic-extraction}, the correction at the target
is $2r_j(\mathbf n)$, so every residual vanishes.

For completeness choose a power of two $b$ exceeding all coordinates
of a given zero, and use its digit code. Since the zero is not the
zero vector, $g>0$. The geometric estimate using $B>b$ gives
$g\le(b-1)\sum_{e=1}^M B^e<bB^M$. The same digit calculations,
now in the forward direction, prove both carry tests. The choice of
$b$ can be made arbitrarily large.
\end{proof}

The central-binomial carry lemma of Jones~\cite{jones1982}, provided
as \texttt{Jones1982.lemma\_2\_16} in the pinned
\texttt{Carries.lean}, states that for $N=2^a$ and $0\le S,T<N$,
\begin{equation}
 \tau_2(S,T)=0
 \quad\longleftrightarrow\quad
 N^2\mid\binom{2R}{R},
 \qquad R=S(N^2-N)+(T+1)(N^2-1).
 \label{eq:carry-binomial}
\end{equation}
This theorem uses precisely the source expression
\eqref{eq:nine-source}. It needs no quartic index or base-five
normalization. The source construction is now complete.

\subsection{The full nine-variable polynomial}

The nine nonnegative unknowns are
\[
 (\epsilon,g_0,h_0,s_0,w_0,\phi,i_0,j_0,n)\in\N^9.
\]
Put $b=\epsilon+2$ and
$g=g_0+1,h=h_0+1,s=s_0+1,w=w_0+1,i=i_0+1,j=j_0+1$.
Use $R,N,G_{\rm cap}$ from \eqref{eq:nine-source} and define
\begin{align}
 U&=N^2w,&Y&=N^2s,&P&=2(RY)^2U,\nonumber\\
 K&=R+1+h(P-1),&A&=RY(U+1),&B_1&=2R+1,\nonumber\\
 W_1&=bw,&C&=B_1+W_1+\phi,&J&=4A-5,\nonumber\\
 Q&=3W_1C-2(W_1^2-1),&D&=(A^2-1)C^2+1,\nonumber\\
 E&=iJDC^2,&F&=(A^2-1)E^2+1,\nonumber\\
 G_p&=A+F(CD+1-A),&H&=B_1+jC,&I&=(G_p^2-1)H^2+1,\nonumber\\
 a_0&=(P^2-1)K^2+1,&a_1&=DFI,&u&=JF,\nonumber\\
 v&=QF+J(H-C),&&&\nonumber\\[-1mm]
 m&=K^2G_{\rm cap}-4(C-KY)^2G_{\rm cap}-8gK^2.&&&
 \label{eq:jones-ratio}
\end{align}
Here $m$ is a register; the ninth unknown is $n$. Next set
\begin{align}
 \sigma&=Y(U+1),&r&=W_1+\phi,\nonumber\\
 M_D&=AC+\sigma(r+1)^2+1,&M_F&=AE+1,\nonumber\\
 T_{CH}&=CH+r^2,&
 M_I&=F(DT_{CH}+2H^2+A^2)-A^2+1,\nonumber\\
 V_0&=PK+1,&V_1&=M_DM_FM_I+1,\nonumber\\
 M_{\rm div}&=Q^2+(H-C)^2+1,&
 \Omega&=u(2m-1),\nonumber\\
 \Theta&=u(n+M_{\rm div})+v
          -\Omega(2M_{\rm div}+V_0V_1).&&
 \label{eq:majorants}
\end{align}
The required equation is
\begin{equation}
 \boxed{F_9=
   \bigl(\Theta^2-\Omega^2(a_0+V_0^2a_1)\bigr)^2
    -4\Omega^4V_0^2a_0a_1=0.}
 \label{eq:F9}
\end{equation}
Equations \eqref{eq:weights}--\eqref{eq:F9}, together with the four
native ports, specify a literal integer-coefficient polynomial. The
JSON arithmetic circuit is its expanded straight-line definition.
No square roots, tests, divisions, or variable exponents occur in it.

\subsection{Signed bounds before positive source recovery}

It is essential not to assume $R>0$ when applying the final polynomial.
For a freely chosen large $g$, $E_*(B)$ can make $R$ negative. The
following proof makes the order of implications explicit.

By Lemma~\ref{lem:source-size}, $R\ne0$ for every assignment of
the eight outer natural coordinates. Since $U,Y\ge1$,
\[
 |A|\ge2,\quad P\ge2R^2\ge2,\quad
 K\ge R+P\ge2R^2+R\ge1.
\]
Consequently $J\ne0$, $AJ>0$, $D\ge1$, $AE\ge0$, and $F\ge1$.
The completed-square bounds from package~37 apply with no sign
restriction on $R$. For completeness, here are their short proofs.
Because $A=R\sigma$ and $C=2R+r+1$,
\begin{align*}
 M_D-AC&=\sigma(r+1)^2+1>0,\\
 M_D+AC&=\sigma\bigl(4R^2+2R(r+1)+(r+1)^2\bigr)+1>0.
\end{align*}
The quadratic expression is nonnegative; twice it is
$(2R+r+1)^2+4R^2+(r+1)^2$.
Thus $M_D\ge|AC|+1$ and
$D=(AC)^2-C^2+1\le M_D^2$.
Similarly $M_F=AE+1\ge1$ and $F\le M_F^2$.

Since $H=(j+1)C-r$,
\[
 T_{CH}-CH=r^2\ge0,\qquad
 T_{CH}+CH=C^2+(C-r)^2+2jC^2\ge0.
\]
It follows that $T_{CH}\ge|CH|$.
Using $G_p=FDC+F-A(F-1)$, the integer inequalities
$|H|\le H^2$ and $|AH|\le A^2+H^2$ imply
\[
 |G_pH|+|H|+1
 \le FDT_{CH}+2FH^2+(F-1)A^2+1=M_I.
\]
Thus $M_I\ge1$ and $\sqrt{|I|}\le M_I$, even when the candidate
$I$ is negative. Altogether
\begin{equation}
 V_0,V_1\ge2,\qquad
 \sqrt{|a_t|}\le V_t-1\quad(t=0,1).
 \label{eq:root-bounds}
\end{equation}
Indeed $1\le a_0\le(PK)^2$ and
$\sqrt{|a_1|}\le M_DM_FM_I$. Also $u\ne0$ and
\begin{equation}
 \left|\frac vu\right|
 =\left|\frac QJ+\frac{H-C}{F}\right|
 \le|Q|+|H-C|\le M_{\rm div}-1.
 \label{eq:quotient-bound}
\end{equation}
All these statements precede square conditions, divisibility, and $m>0$.

\begin{lemma}[Signed-divisor finalizer]
\label{lem:finalizer}
Let integers $a_0,a_1,u,v,m,M,V_0,V_1$ satisfy $u\ne0$, $M\ge1$,
$|v/u|\le M-1$, $V_0,V_1\ge2$, and
$\sqrt{|a_t|}\le V_t-1$. Define $\Omega$, $\Theta$, and $F_9$ as
above, with $M_{\rm div}=M$. Then
\[
 \exists n\in\N\ F_9=0
 \quad\longleftrightarrow\quad
 a_0,a_1\text{ are integer squares},\quad u\mid v,\quad m>0.
\]
\end{lemma}
\begin{proof}
The odd integer $2m-1$ is nonzero, so $\Omega\ne0$. Choose complex
roots $r_t^2=a_t$. Direct multiplication factors $F_9$ into
\[
 \prod_{\varepsilon,\delta\in\{-1,1\}}
       \bigl(\Theta-\Omega(\varepsilon r_0+\delta V_0r_1)\bigr).
\]
A zero gives one vanishing factor. Absorb its signs in the chosen
roots, and put $s=\Theta/\Omega=r_0+V_0r_1\in\Q$.
If $s\ne0$, the identity
$2sr_0=s^2+a_0-V_0^2a_1$ makes both roots rational. A rational
number with integer square is an integer. If $s=0$, the bounds give
$|r_1|=|r_0|/V_0<1$, hence the integer $|a_1|<1$; then
$a_1=a_0=0$. This includes negative candidate radicands and coincident
square classes, without presuming independent field automorphisms.

Set $Z=2M+V_0V_1+r_0+V_0r_1$. The root bounds give
$Z\ge2M+1>0$. Dividing the vanishing factor by the possibly negative
$u$ gives
\[
 n+M+v/u=(2m-1)Z.
\]
Its left side is at least $n+1>0$, so $m>0$. Its right side is
integral, so $u\mid v$. Conversely, if these three conditions hold,
choose integer roots and define
$n=(2m-1)Z-M-v/u$. This is an integer at least $2$, and makes
the corresponding factor vanish.
\end{proof}

\subsection{Soundness and completeness of the nine-variable form}

The parameterized ratio theorems in
\texttt{NineRatio.lean} have the following precise interface.
For $R,N\ge8$, $0<b\le\min\{R,N\}$, and $G>0$, the conditions
\[
 a_0,a_1\text{ squares},\qquad JF\mid QF+J(H-C),\qquad m>0
\]
with $h,s,w,i,j>0$ and $\phi\ge0$ imply
\begin{equation}
 N^2\mid\binom{2R}{R},\qquad b\text{ is a power of two},
 \qquad8g<G.
 \label{eq:ratio-sound}
\end{equation}
Conversely, the first two conditions in \eqref{eq:ratio-sound}
together with $32g<G$ produce positive $h,s,w,i,j$ and
nonnegative $\phi$ satisfying the three displayed ratio conditions.
These are \texttt{Jones1982.Nine.ratio\_sound} and
\texttt{ratio\_complete}; the formulas match \eqref{eq:jones-ratio}
with $G=G_{\rm cap}$. They quantify arbitrary admissible $R,N,b,G$
and do not require Jones's original source packing.

Suppose $F_9=0$. The unconditional bounds and
Lemma~\ref{lem:finalizer} first give $m>0$. As $K\ge1$ and the
subtracted square term in $m$ is nonnegative, this already implies
\[
 8g<G_{\rm cap},\qquad g<4bB^M.
\]
Only now apply Lemma~\ref{lem:source-size}: it supplies positive
$R,N$ and every size hypothesis of the ratio theorem. Apply
\eqref{eq:ratio-sound}, then \eqref{eq:carry-binomial}, then
Lemma~\ref{lem:carry-decoder}. This gives a zero of the 22 parent
residuals, and hence a blank K215 halt by
\eqref{eq:parent-calibration}.

Conversely suppose the parent residuals have a natural zero. Choose
$b$ and its positive code $g<bB^M$ by
Lemma~\ref{lem:carry-decoder}. Its carry condition and
\eqref{eq:carry-binomial} give the binomial divisibility; all ratio
size conditions hold. Since $32g<32bB^M=G_{\rm cap}$,
\texttt{ratio\_complete} supplies the six remaining outer witnesses.
Lemma~\ref{lem:finalizer} supplies $n\in\N$. Translating the positive
outer witnesses down by one gives a zero of $F_9$. We have proved
\[
 \exists\N^9\ F_9=0
 \quad\longleftrightarrow\quad
 \exists\N^{31}\ \bigwedge_jr_j=0,
\]
and thus the first line of Theorem~\ref{thm:main}.

For strictly positive coordinates use
\[
 (\epsilon_+,g,h,s,w,\phi_+,i,j,n_+)\in\Zpos^9,
 \quad b=\epsilon_++1,\quad\phi=\phi_+-1,\quad n=n_+-1.
\]
This coordinate bijection gives $F_9^+$, preserves exact degree,
and saves four operations in the delivered evaluator. It does not
extend the result to nine unrestricted integers.

\subsection{Exact degree, with cancellations resolved}

Treat fixed coefficients as degree zero. The actual degree cannot
be read from the original nested expression by blindly taking
maxima: its highest apparent terms cancel. The following exact
identities expose the surviving leading term. Put
$t_0=\sigma(r+1)^2$ and
\begin{align*}
 \delta_D&=M_D^2-D=2ACt_0+t_0^2+2AC+2t_0+C^2,\\
 \delta_F&=M_F^2-F=2AE+E^2,\\
 W_-&=M_I-G_pH\\
 &=FDr^2+2FH^2+(F-1)A^2+1-FH+A(F-1)H,\\
 \delta_I&=M_I^2-I=W_-(M_I+G_pH)+H^2-1,\\
 \delta_V&=V_1^2-DFI\\
 &=\delta_DM_F^2M_I^2+D\delta_FM_I^2+DF\delta_I
                      +2M_DM_FM_I+1.
\end{align*}
With $E_0=u(n+M_{\rm div})+v-2\Omega M_{\rm div}$, the inner
square in \eqref{eq:F9} is exactly
\begin{equation}
 E_0^2-2E_0\Omega V_0V_1+\Omega^2V_0^2\delta_V-\Omega^2a_0.
 \label{eq:degree-inner}
\end{equation}
Some decisive total degrees are as follows:
\[
\begin{array}{c|r@{\qquad}c|r}
 N&142270&R&426808\\
 P&1707239&K&1707240\\
 A&995890&C&426808\\
 D&2845396&F&11381586\\
 I&30161198&a_1&44388180\\
 V_0&3414479&V_1&22194090\\
 \Omega&16362637&\delta_D&1991784\\
 \delta_F&9389806&\delta_I&29307585\\
 \delta_V&43534568&E_0&17216257
\end{array}
\]
The first three terms defining $\delta_V$ have degrees
$43\,534\,568$, $42\,396\,400$, and $43\,534\,567$.
The first uniquely dominates. The third term in
\eqref{eq:degree-inner} uniquely dominates, at degree $83\,088\,800$.
The final subtracted term of \eqref{eq:F9} has strictly smaller degree
than twice this value. Hence
\begin{equation}
 \boxed{\deg F_9=166\,177\,600.}
 \label{eq:degree9}
\end{equation}
For a nonvanishing certificate, substitute all nine natural inputs by
$t$. The coefficient of $t^{166177600}$ is exactly
\[
 \boxed{2^{196}\kappa^{83088652}>0.}
\]
The verifier tracks at every gate a total-degree upper bound and the
exact coefficient of the corresponding diagonal power. All sixteen
derived native coefficients remain independent symbolic leaves; none
is assigned zero or presumed generic. The final coefficient is
independent of them. This proves the claimed degree for the actual
materialized tuple as well as for every compatible specialization.

The production circuits use sparse Horner evaluation of $E_*$,
shared repeated squaring, and a divide-and-conquer geometric-sum
identity. Their literal gate counts are 868 for $\N^9$ and 864 for
$\Zpos^9$. The separate cancellation certificate uses 917 operations;
it is a proof artifact, not the claimed fastest evaluator.
